\documentclass[11pt,a4paper]{article}
\usepackage{amsmath,amssymb}
\usepackage{listings}
\usepackage{hyperref}
\usepackage[margin=1in]{geometry}
\usepackage{fancyhdr}

\pagestyle{fancy}
\fancyhf{}
\fancyhead[L]{Module 11: Data Structures Fundamentals}
\fancyhead[R]{C Programming Course}
\fancyfoot[C]{\thepage}

\lstset{
  language=C,
  basicstyle=\ttfamily\small,
  keywordstyle=\bfseries,
  commentstyle=\itshape,
  numbers=left,
  numberstyle=\tiny,
  frame=single,
  breaklines=true,
  tabsize=4
}

\title{Module 11: Data Structures Fundamentals}
\author{C Programming Course}
\date{}

\begin{document}
\maketitle
\thispagestyle{fancy}

\section{Linked Lists}
\subsection{Singly Linked List}
\begin{lstlisting}
#include <stdio.h>
#include <stdlib.h>

typedef struct Node {
    int data;
    struct Node *next;
} Node;

Node *create_node(int data) {
    Node *n = malloc(sizeof(Node));
    if (n == NULL) return NULL;
    n->data = data;
    n->next = NULL;
    return n;
}

void print_list(const Node *head) {
    while (head != NULL) {
        printf("%d -> ", head->data);
        head = head->next;
    }
    printf("NULL\n");
}

void free_list(Node *head) {
    while (head != NULL) {
        Node *tmp = head;
        head = head->next;
        free(tmp);
    }
}
\end{lstlisting}

\subsection{Doubly and Circular Lists}
A doubly linked list adds a \texttt{prev} pointer. A circular list connects the
last node back to the first.

\section{Stacks and Queues}
\begin{itemize}
  \item \textbf{Stack} -- LIFO: push and pop from the top.
  \item \textbf{Queue} -- FIFO: enqueue at the back, dequeue from the front.
\end{itemize}

Both can be implemented with arrays or linked lists.

\section{Hash Tables}
A hash table maps keys to values using a hash function. Collisions are resolved
with chaining (linked lists) or open addressing.

\begin{lstlisting}
#define TABLE_SIZE 64

unsigned int hash(const char *key) {
    unsigned int h = 0;
    while (*key) {
        h = h * 31 + (unsigned char)(*key);
        key++;
    }
    return h % TABLE_SIZE;
}
\end{lstlisting}

\section{Complexity Analysis}
\begin{center}
\begin{tabular}{lcc}
\hline
\textbf{Operation} & \textbf{Array} & \textbf{Linked List} \\
\hline
Access by index & $O(1)$ & $O(n)$ \\
Insert at head  & $O(n)$ & $O(1)$ \\
Search          & $O(n)$ & $O(n)$ \\
\hline
\end{tabular}
\end{center}

\end{document}
